\chapter{Observability and Consistency of VIO}
\label{ch:observability}

\begin{objectives}
\begin{itemize}
  \item Define observability for nonlinear and linearized systems.
  \item Prove that VIO has four unobservable directions and write them in the project's error coordinates.
  \item Understand why the standard EKF gains spurious information about yaw and becomes overconfident.
  \item Know the remedies (FEJ, observability constraints, invariant filtering) and the degenerate motions that remove
        further information.
\end{itemize}
\end{objectives}

% ------------------------------------------------------------------------------
\section{Observability}

A system is \emph{observable} if its state can be determined from its inputs and outputs. For the nonlinear systems of navigation
the relevant notion is local: a state is \emph{locally weakly observable} if it can be distinguished from all its neighbors
using the outputs over some time interval. For an estimator the linearized version is the operative one.

\begin{definition}[Observability matrix]
For the linearized time-varying system $\delta x_{k+1} = \Phi_k\delta x_k$, $r_k = H_k\delta x_k$, the observability matrix over steps $1\dots k$ is
\begin{equation}
  \mathcal{O} = \begin{bmatrix} H_1 \\ H_2\Phi_1 \\ H_3\Phi_2\Phi_1 \\ \vdots \\ H_k\Phi_{k-1}\cdots\Phi_1 \end{bmatrix}.
\end{equation}
An initial error $\delta x_1 \in \nullsp(\mathcal{O})$ produces no change in any measurement, so it cannot be detected.
$\nullsp(\mathcal{O})$ is the \emph{unobservable subspace}.
\end{definition}

A Kalman filter can only gain information in directions outside $\nullsp(\mathcal{O})$. Along unobservable directions, its
covariance may grow (from process noise) but must never shrink. This is the property the standard EKF violates.

% ------------------------------------------------------------------------------
\section{The four unobservable directions of VIO}

\begin{theorem}[Unobservable directions of VIO]
\label{thm:vio-unobservable}
For monocular VIO with point features, applying the same global translation $t$ and the same rotation $R_z(\alpha)$ about the
gravity direction to all positions, velocities, attitudes, clones and features leaves every IMU and camera measurement
unchanged. The state is therefore unobservable along these four directions.
\end{theorem}
\begin{proof}
Under $p \mapsto R_zp + t$, $v \mapsto R_zv$, $R_{wb} \mapsto R_zR_{wb}$, $p_f \mapsto R_zp_f + t$, with biases unchanged:
\begin{itemize}
  \item the gyroscope measures $\sx{\omega_b} = R_{wb}\T\dot R_{wb}$, which is invariant because $R_z$ is constant;
  \item the accelerometer measures $R_{wb}\T R_z\T(R_za_w - g_w) = R_{wb}\T(a_w - R_z\T g_w) = R_{wb}\T(a_w - g_w)$, because a rotation about
        gravity leaves $g_w$ unchanged;
  \item the camera sees $R_{bc}\T\big(R_{wb}\T R_z\T(R_zp_f + t - R_zp - t) - p_{bc}\big) = R_{bc}\T\big(R_{wb}\T(p_f - p) - p_{bc}\big)$.
\end{itemize}
All measurements are identical, so the transformed and original trajectories cannot be distinguished.
\end{proof}

Rotations about horizontal axes are \emph{not} in this family, because they would rotate gravity in the accelerometer equation;
roll and pitch are observable. Scaling is not either, because it would scale the accelerations the accelerometer measures;
scale is observable as long as the vehicle accelerates.

Linearizing the transformation gives the unobservable directions in the project's error coordinates. With $e_z$ the vertical
unit vector, a small yaw rotation $\alpha$ and a small translation $t$ change the state by (verified numerically):
\begin{equation}
  \underbrace{\begin{bmatrix} \delta p \\ \delta v \\ \delta\theta \\ \delta b_a \\ \delta b_g \\ \delta p_i \\ \delta\theta_i \\ \delta p_f \end{bmatrix}}_{\text{error state}}
  = \underbrace{\begin{bmatrix}
    I_3 & -\sx{\hat p}\,e_z \\
    0 & -\sx{\hat v}\,e_z \\
    0 & \hat R\T e_z \\
    0 & 0 \\
    0 & 0 \\
    I_3 & -\sx{\hat p_i}\,e_z \\
    0 & \hat R_i\T e_z \\
    I_3 & -\sx{\hat p_f}\,e_z
  \end{bmatrix}}_{N}
  \begin{bmatrix} t \\ \alpha \end{bmatrix}.
  \label{eq:nullspace-N}
\end{equation}
The last row applies only to formulations that keep features in the state. The attitude entry is $\hat R\T e_z$ because the error is
local: a rotation about the world vertical, expressed in body axes. Every column of $N$ \emph{depends on the estimate}: the yaw
column involves $\hat p$, $\hat v$, $\hat R$ and the clone and feature estimates. That dependence is the root of the problem.

% ------------------------------------------------------------------------------
\section{Why the standard EKF becomes overconfident}

If every Jacobian were evaluated at the true state, the linearized system would satisfy $\mathcal{O}N = 0$ at all times, and the
filter would never gain information along $N$. The EKF instead evaluates each Jacobian at the best estimate available when it is
computed. The transition matrix $\Phi_k$ is evaluated at the propagated estimate of step $k$, while the next measurement Jacobian is
evaluated at an estimate of the same state that has since been corrected by an update. The same physical state enters the
observability matrix linearized at two different points, and the products $H_{k}\Phi_{k-1}\cdots\Phi_1 N$ no longer cancel.

The translation columns survive this, because all Jacobians depend only on \emph{differences} of positions. The yaw column does
not: with estimate-dependent Jacobians, $\mathcal{O}N_{\text{yaw}} \neq 0$, and the unobservable subspace of the EKF's linearized model has
dimension three instead of four~\cite{huang2010observability, li2013msckf2, hesch2014consistency}. The filter then extracts
\emph{spurious information} about the global yaw from measurements that contain none. Its yaw variance shrinks when it should not, it
becomes overconfident in yaw, and because yaw errors couple into position through the trajectory, its position covariance becomes
too small as well. On long trajectories the NEES of yaw and position grows steadily while the NIS can look perfectly healthy: the
inconsistency lives in directions the measurements do not see (\cref{tab:consistency-diagnosis}).

\begin{keyidea}[an inconsistency that no unit test catches]
EKF yaw inconsistency is not a bug in the usual sense: every Jacobian can pass its finite-difference test. It is a property of
linearizing at estimates that change over time. Milestone 10b therefore expects it, and asks for it to be measured and documented
rather than fixed.
\end{keyidea}

% ------------------------------------------------------------------------------
\section{Remedies}

\paragraph{First-Estimates Jacobians (FEJ).} Evaluate every Jacobian that involves a state at the \emph{first} estimate of that
state, the value it had when it entered the filter, rather than at the latest one~\cite{huang2008fej, li2013msckf2}. Each state then
appears in the observability matrix linearized at a single point, the products telescope, and $\mathcal{O}N = 0$ holds again with $N$
built from first estimates. In the MSCKF this means using the propagated (pre-update) estimates of position and velocity at both
ends of each propagation interval in $\Phi$, and the first estimates of the clone poses and of the triangulated feature in the
measurement Jacobians. The price is a somewhat larger linearization error; in exchange the filter is markedly more consistent and,
in practice, more accurate. FEJ is the project's stretch goal.

\paragraph{Observability-constrained EKF (OC-EKF).} Modify each Jacobian minimally, in the least-squares sense, so that the correct
unobservable subspace is enforced explicitly~\cite{hesch2014consistency}.

\paragraph{Invariant EKF.} Define the error on the extended pose group $\mathrm{SE}_2(3)$ with a right-invariant error. The error
dynamics and the unobservable directions then become independent of the estimate, and consistency with respect to these directions
holds by construction~\cite{barrau2017invariant}.

\paragraph{Robocentric formulations.} Express the state relative to the current body frame, so that the unobservable global
quantities are not estimated at all.

% ------------------------------------------------------------------------------
\section{Degenerate motions}

\Cref{thm:vio-unobservable} assumes \emph{general} motion. Some motions remove further directions~\cite{yang2019degenerate}.
\begin{center}\small
\begin{tabular}{@{}l l@{}}
\toprule
Motion & What becomes unobservable \\
\midrule
constant velocity (no acceleration) & metric scale \\
pure rotation (no translation) & feature depths and translation scale \\
no rotation & accelerometer bias inseparable from the gravity direction (tilt) \\
at rest & both of the above, plus everything that needs parallax \\
\bottomrule
\end{tabular}
\end{center}
The filter does not fail during such motions, provided it is consistent: its covariance simply stops shrinking along the
affected directions until informative motion returns. An inconsistent filter, by contrast, may ``learn'' these quantities from
noise. When a vehicle is known to be stationary, a \emph{zero-velocity update} (the pseudo-measurement $v = 0$) supplies real
information and stops drift, which is why many systems detect stillness explicitly.

\begin{takeaways}
\begin{itemize}
  \item The unobservable subspace is the nullspace of the observability matrix; a filter must not gain information along it.
  \item VIO has four unobservable directions: global translation and rotation about gravity. Roll, pitch and scale are observable
        under general motion.
  \item In local error coordinates the yaw direction is $\hat R\T e_z$ for the attitude and $-\sx{\hat p}e_z$ for positions: it depends on
        the estimate.
  \item Linearizing the same state at different estimates makes yaw spuriously observable, so the EKF becomes overconfident in yaw and
        position.
  \item FEJ, OC-EKF and invariant filtering restore consistency; degenerate motions remove further directions.
\end{itemize}
\end{takeaways}

\begin{exercises}
  \item Verify \eqref{eq:nullspace-N} numerically: apply a small yaw rotation and translation to a random state and compute
        $x'\bminus x$.
  \item Show that the translation columns of $N$ remain in the nullspace of the EKF's observability matrix even when Jacobians are
        evaluated at changing estimates.
  \item For a planar robot with odometry and range-bearing measurements to a landmark, show that the EKF gains spurious
        heading information. (This is the classic EKF-SLAM inconsistency result.)
  \item Implement a zero-velocity update as an EKF measurement and derive its $H$. When is it safe to apply?
  \item Read~\cite{li2013msckf2} and list which Jacobians must be evaluated at first estimates in the MSCKF. Sketch how you would store
        the first estimates in the project's data structures without allocation.
\end{exercises}
